\de{ĐỀ THI HỌC KỲ I NĂM HỌC 2022-2023}{THPT Bình Phú}



\begin{bt}%[0D1Y3-2]%[Dự án đề kiểm tra HKI NH22-23-VUNGocHao]%[THPT Bình Phú]
	Cho các tập hợp $A=\{2 ; 4 ; 6 ; 8\}, B=\{2 ; 3 ; 4 ; 5 ; 6 ; 7\},\,C=(2 ; 4)$.
	Tìm $A \cap B$, $A\cap C$, $A \cup B$ và $A\setminus B$.
	\loigiai{
	Ta có \begin{itemize}
		\item $A \cap B=\{2;4;6\}$.
		\item $A\cap C=\varnothing$.
		\item $A\cup B=\{2 ; 3 ; 4 ; 5 ; 6 ; 7;8\}$.
		\item $A\setminus B=\{8\}$.
	\end{itemize}
	}
\end{bt} 
%%%%%%%%%%%============Câu 2
\begin{bt}%[0H4B2-1]%[Dự án đề kiểm tra HKI NH22-23-VUNGocHao]%[THPT Bình Phú]
	\begin{enumerate}
		\item Tính $\cos \alpha$, biết rằng $\sin \alpha=\dfrac{2}{3}$ với $90^{\circ}<\alpha<180^{\circ}$.
		\item Cho $\triangle ABC$ có $AB=6,\,AC=8,\,\widehat{A}=60^{\circ}$. Tính $BC$, diện tích và bán kính đường tròn ngoại tiếp tam giác $ABC$.
	\end{enumerate}
	\loigiai{
	\begin{enumerate}
		\item Ta có $\sin^2 \alpha + \cos^2 \alpha=1 \Rightarrow \cos^2 \alpha =1-\sin^2\alpha = 1-\dfrac{4}{9}=\dfrac{5}{9}\Rightarrow \hoac {&\cos\alpha=-\dfrac{\sqrt{5}}{3}\\
			&\cos\alpha=\dfrac{\sqrt{5}}{3}.}$\\
		Do $90^{\circ}<\alpha<180^{\circ}$ nên chọn $\cos\alpha=-\dfrac{\sqrt{5}}{3}$.
		\item 	
		\begin{itemize}
			\item 	Áp dụng định lý cô-sin trong $\triangle ABC$ ta có\\ $BC^2=AB^2+AC^2-2AB\cdot AC\cdot \cos 60^{\circ}=36+64-2\cdot6\cdot8\cdot\dfrac{1}{2}=52$\\
			$\Rightarrow BC = 2\sqrt{13}.$
			\item Diện tích $\triangle ABC$ 
			$$S_{\triangle ABC}=\dfrac{1}{2}\cdot AB\cdot AC \cdot\sin A= \dfrac{1}{2}\cdot6\cdot 8 \cdot\sin 60^{\circ}=12\sqrt{3}.$$
			\item $S=\dfrac{AB\cdot AC\cdot BC}{4R} \Rightarrow R= \dfrac{6\cdot 8\cdot 2\sqrt{13}}{4\cdot 12\sqrt{13}}=2.$
		\end{itemize}
	
		
	\end{enumerate}
		

	}
\end{bt}
%%%%%%%%%%%============Câu 3
\begin{bt}%[0D3K2-1]%[Dự án đề kiểm tra HKI NH22-23-VUNGocHao]%[THPT Bình Phú]
	Một cửa hàng điện lạnh dự định kinh doanh hai loại máy điều hòa, điều hòa hai chiều và điều hòa một chiều với số vốn ban đầu không vượt quá $1{,}2$ tỉ đồng ($1\,200$ triệu đồng).
	\begin{center}
		\begin{tabular}{|l|l|l|}
			\hline & \textbf{Điều hòa 2 chiều} & \textbf{Điều hòa 1 chiều} \\
			\hline \textbf{Giá mua vào} &$  20 $ triệu đồng/1 máy & $ 10 $ triệu đồng/1 máy \\
			\hline \textbf{Lợi nhuận dự kiến} & $ 3{,}5 $ triệu đồng/1 máy&$2$ triệu đồng/1 máy \\
			\hline
		\end{tabular}
	\end{center}
	Cửa hàng ước tính rằng tổng nhu cầu của thị trường sẽ không vượt quá $100$ máy cả hai loại. Nếu là chủ cửa hàng thì em cần đầu tư kinh doanh mỗi loại bao nhiêu máy đề lợi nhuận thu được là lớn nhất.
	\loigiai{
		Gọi $x$ là số máy điều hòa hai chiều ($0< x\leq 100$).\\
		Số máy điều hòa một chiều là $100-x$.\\
		Lợi nhuận thu được là $L=3{,}5x+2\cdot(100-x)=1{,}5x+200$.\\
		Do số vốn ban đầu không vượt quá $1{,}2$ tỉ đồng nên ta có: $$
		20x+10\cdot (100-x)\leq 1200\Leftrightarrow x\leq 20.
		$$
		Do \[\begin{aligned}
			&x\leq 20\\\Leftrightarrow& 1{,}5x\leq 30\\\Leftrightarrow&1{,}5x+200\leq 230\\\Leftrightarrow&L\leq 230.
		\end{aligned}\]
		Suy ra $L_{\text{max}}=230\Leftrightarrow x=20$.\\
		Vậy để thu được lợi nhuận lớn nhất thì cần kinh danh $20$ máy điều hòa hai chiều và $80$ máy điều hòa một chiều.
	}
\end{bt} 


	\begin{bt}%[0T3K2-5]%[Dự án đề kiểm tra HKII-NH22-23, Dương Phước Sang]%[THPT Bình Phú-TP.HCM]
		\immini{
			\begin{enumerate}
				\item Vẽ parabol $(P)\colon y=x^2-4 x+3$.
				\item Khi một quả bóng được đá lên, nó sẽ đạt đến độ cao nào đó rồi rơi xuống. Biết rằng quỹ đạo của quả bóng là một cung parabol trong mặt phẳng với hệ tọa độ $Oxy$, trong đó $x$ là thời gian (tính bằng giây) kể từ khi quả bóng được đá lên và $y$ là độ cao (tính bằng mét) của quả bóng. Giả thiết rằng quả bóng được đá từ một nóc nhà cao $3 \mathrm{~m}$. Sau đó $1$ giây, quả bóng đạt độ cao $6 \mathrm{~m}$ và $4$ giây sau khi đá lên, nó ở độ cao bằng với độ cao từ vị trí xuất phát (xem hình vẽ bên). Xác định độ cao lớn nhất của quả bóng.
			\end{enumerate}}
			{\vspace{-0.5cm}
				\begin{tikzpicture}[>=stealth,x=1cm,y=1cm,scale=.7,font=\footnotesize]
					\def\a{-1}
					\def\b{4}
					\def\c{3}
					\path
					(0,3) coordinate (A)
					(4.65,0) coordinate (B)
					(0,0)coordinate (O);
					\draw[->] (-1,0) -- (5.5,0) node[below] {\scriptsize $x$};
					\draw[->] (0,-1) -- (0,7.5) node[left] {\scriptsize $y$};
					\draw[thick,samples=150,smooth,domain=0:4.65] plot(\x,{\a*(\x)^2+(\b)*\x+(\c)});
					\draw 
					(0,3)node[below right]{$3$}
					(0,6)node[left]{$6$}
					(4,0)node[below]{$4$}
					(1,0)node[below]{$1$};
					\draw[dashed] (0,3)-|(4,0) (0,6)-|(1,0);
					\foreach \x/\g in {A/180,B/-90,O/-135}\fill (\x) circle (.05) +(\g:.4) node{$\x$};
			\end{tikzpicture}}
		\dapso{b) $7$ m.  }
		\loigiai{
			\begin{enumerate}
				\item Parabol $(P)\colon y=x^2-4 x+3$ có trục đối xứng $x=2$, tọa độ đỉnh $I(2;-1)$.\\
				Bảng biến thiên
				\begin{center}
					\begin{tikzpicture}
						\tkzTabInit[nocadre=true,lgt=1.2,espcl=2.5,deltacl=0.6]
						{$x$/0.6,$y$/2}
						{$-\infty$,$2$,$+\infty$}
						\tkzTabVar{+/$+\infty$,-/$-1$,+/$+\infty$}
					\end{tikzpicture}
				\end{center}
				Đồ thị hàm số $y=x^2-4 x+3$ như sau
				\begin{center}
					\begin{tikzpicture}[scale=0.7, font=\footnotesize, samples=200, >=stealth]
						\draw[->] (-1,0)--(5,0) node[below]{$x$};
						\draw[->] (0,-1.3)--(0,4) node[left]{$y$};
						\node at (0,0) [below left=-2pt]{$O$};
						\draw plot[domain=-0.2:4.2](\x,{(\x)^2-4*\x+3});
						\fill
						(2,-1) circle(1.07pt)
						(1,0) circle(1.07pt) node[below left=-3pt]{$1$}
						(3,0) circle(1.07pt) node[below right=-3pt]{$3$}
						(3,0) circle(1.07pt);
						\draw[dashed]
						(2,0) node[above]{$2$} |-(0,-1)node[left]{$-1$};
					\end{tikzpicture}
				\end{center}
				\item Quỹ đạo bay của quả bóng là đồ thị hàm số $y=h(x)=ax^2+bx+c$. Theo giả thiết ta có
				$$\heva{&h(0)=3\\&h(1)=6\\&h(4)=3} 
				\Leftrightarrow \heva{&c=3\\&a+b+c=6\\&16a+4b+c=3}
				\Leftrightarrow \heva{&a=-1\\&b=4\\&c=3.}$$
				Như vậy $h(x)=-x^2+4x+3$ và độ cao lớn nhất của quá bóng đạt tại đỉnh của quỹ đạo, đó là $h(2)=7 \mathrm{\,m}$.
			\end{enumerate}
		}
	\end{bt} 
	
	\begin{bt}%[0T5K4-1]%[Dự án đề kiểm tra HKII-NH22-23, Dương Phước Sang]%[THPT Bình Phú-TP.HCM]
		Cho tam giác $ABC$ vuông tại $B$ có $M$, $N$ lần lượt là trung điểm các cạnh $AB$ và $AC$.
		\begin{enumerate}
			\item Chứng minh rằng $\overrightarrow{BC}=2\overrightarrow{AN}-2 \overrightarrow{AM}$.
			\item Gọi $H$ là điểm trên cạnh $AC$ sao cho $\overrightarrow{AH}\cdot \overrightarrow{AC}=AB^2$. Chứng minh $BH\perp AC$.
		\end{enumerate}
		%\dapso{}
		\loigiai{
			\immini{
				\begin{enumerate}
					\item[a)] Do $M$ và $N$ lần lượt là trung điểm các cạnh $AB$ và $AC$ nên 
					$$\heva{&\vec{AB}=2\vec{AM}\\&\vec{AC}=2\vec{AN}.}$$
					Từ đó $\vec{BC}=\vec{AC}-\vec{AB}=2\vec{AN}-2\vec{AM}$.
			\end{enumerate}}
			{\begin{tikzpicture}[scale=1, font=\footnotesize, line join=round, line cap=round]
					\coordinate (A) at (0,0);
					\coordinate (B) at (60:2);
					\coordinate (C) at (4,0);
					\coordinate (M) at ($(A)!0.5!(B)$);
					\coordinate (N) at ($(A)!0.5!(C)$);
					\draw (A)--(B)--(C)--cycle;
					\node at (A) [below left=-3pt]{$A$};
					\node at (B) [above=-1pt]{$B$};
					\node at (C) [below right=-3pt]{$C$};
					\node at (M) [above left=-3pt]{$M$};
					\node at (N) [below=-1pt]{$N$};
					\fill (M) circle (1pt)  (N) circle (1pt);
			\end{tikzpicture}}
			\begin{enumerate}
				\item[b)] Ta có 
				\allowdisplaybreaks
				$\begin{aligned}[t]
					\overrightarrow{AH}\cdot \overrightarrow{AC}=AB^2
					&\Leftrightarrow \left(\overrightarrow{AB}+\overrightarrow{BH}\right)\cdot \overrightarrow{AC}=\overrightarrow{AB}^2\\
					&\Leftrightarrow \overrightarrow{AB} \cdot \overrightarrow{AC}+\overrightarrow{BH}\cdot \overrightarrow{AC}=\overrightarrow{AB}^2\\
					&\Leftrightarrow \overrightarrow{BH}\cdot \overrightarrow{AC}=\overrightarrow{AB}\left(\overrightarrow{AB}-\overrightarrow{AC}\right)\\
					&\Leftrightarrow \overrightarrow{BH}\cdot \overrightarrow{AC}=\overrightarrow{AB}\cdot \overrightarrow{CB}=0 \quad (\text{vì }AB \perp CB).
				\end{aligned}$\\
				Như vậy $BH \perp AC$.
			\end{enumerate}
		}
	\end{bt} 
	
	\begin{bt}%[0T6K4-1]%[Dự án đề kiểm tra HKII-NH22-23, Dương Phước Sang]%[THPT Bình Phú-TP.HCM]
		Một nhà nghiên cứu ghi lại số tuổi của $20$ bệnh nhân. Kết quả thống kê ở bảng số liệu chi tiết sau đây
		\begin{center}
			\begin{tabular}{|c|c|c|c|c|c|c|c|c|c|}
				\hline $21$ & $17$ & $20$ & $18$ & $20$ & $17$ & $15$ & $13$ & $15$ & $20$ \\
				\hline $15$ & $12$ & $18$ & $17$ & $15$ & $16$ & $21$ & $15$ & $12$ & $18$ \\
				\hline
			\end{tabular}
		\end{center}
		\begin{enumerate}
			\item Lập bảng phân bố tần số. Tính số tuổi trung bình của các bệnh nhân?
			\item Tìm mốt và tứ phân vị của mẫu số liệu.
		\end{enumerate}
		\dapso{a) Số tuổi trung bình: $16{,}75$; b) Mốt: $15$; $Q_1=15,\,Q_2=17,\,Q_3=19$. }
		\loigiai{
			\begin{enumerate}
				\item Bảng phân bố tần suất số tuổi của $20$ bệnh nhân
				\begin{longtable}{|c|c|c|c|c|c|c|c|c|c|}
					\hline
					Số tuổi & $12$ & $13$ & $15$ & $16$ & $17$ & $18$ & $20$ & $21$ & $N$\\
					\hline
					Số bệnh nhân & $2$ & $1$ & $5$ & $1$ & $3$ & $3$ & $3$ & $2$ & $20$\\
					\hline
				\end{longtable}
				Số tuổi trung bình của $20$ bệnh nhân được thống kê là
				$$\overline{x}=\dfrac{12\cdot 2+13+15\cdot 5+16+17\cdot 3+18\cdot 3+20\cdot 3+21\cdot 2}{20}=16{,}75.$$
				\item Mốt của mẫu số liệu là $M_0=15$.\\
				Trung vị của mẫu số liệu là $M_e=17$, do đó $Q_2=17$.\\
				Tứ phân vị thứ nhất $Q_1=15$ và tứ phân vị thứ ba là $Q_3=\dfrac{18+20}{2}=19$.\\
				Các tứ phân vị của mẫu số liệu là $Q_1=15,\,Q_2=17,\,Q_3=19$.
			\end{enumerate}
		}
	\end{bt} 


